\section{Discussion et perspectives}

\begin{frame}
    \frametitle{Discussion: limites actuelles}
    \begin{block}{Portée des résultats}
        \footnotesize
        Le pipeline est un \textbf{PoC in silico}: les signaux chimie/biologie guident l'optimisation, mais ne remplacent pas une validation expérimentale.
    \end{block}

    \begin{block}{Limites techniques principales}
        \footnotesize
        \begin{itemize}
            \item Génération actuellement en \textbf{longueur fixe}.
            \item Qualité du contrôle dépendante de la \textbf{qualité des labels} de training.
            \item Score bio ESM utile pour classer, mais \textbf{pas une vérité terrain} d'activité.
        \end{itemize}
    \end{block}
\end{frame}

\begin{frame}
    \frametitle{Perspectives prioritaires}
    \begin{itemize}
        \item \textbf{Entrée en langage naturel}: LLM pour traduire des objectifs utilisateurs en contraintes structurées.
        \item \textbf{Chimie enrichie}: ajouter des descripteurs moléculaires (ex. dérivés SMILES) en complément des propriétés séquence.
        \item \textbf{Biologiste spécialisé peptides}: modèle target-conditionné entraîné/fine-tuné sur données peptides validées.
        \item \textbf{Généralisation bio}: améliorer la robustesse sur espèces peu documentées ou émergentes.
    \end{itemize}
    \begin{center}
        $\implies$ Meilleure fiabilité des entrées \& résultats.
        \\
        $\implies$ Utilisation par public non-expert.
        
    \end{center}
\end{frame}